\chapter{Massless Particles, Helicity, and Gauge Redundancy}
\label{ch:massless-particles-helicity-gauge-redundancy}

\section{The Massless Orbit}

Work in \(D=4\) with mostly-plus metric. The positive-energy massless orbit is
\[
  \Sigma_0^+
  =
  \{k=(k^0,\vec k):k^2=0,\ k^0=|\vec k|,\ k^0>0\}.
\]
Choose a reference momentum
\[
  q=(\kappa,0,0,\kappa),
  \qquad
  \kappa>0.
\]
For each \(k\in\Sigma_0^+\), choose a standard Lorentz transformation
\[
  L(k)\in SO^+(1,3),
  \qquad
  L(k)q=k.
\]
The little group of \(q\) is
\[
  G_q=\{W\in SO^+(1,3):Wq=q\}.
\]

Let \(J_{\mu\nu}\) denote the Lorentz generators acting on one-particle
states, with the same algebra convention as in the massive-spin chapter. The
three generators
\[
  J_3=J_{12},
  \qquad
  N_1=J_{01}+J_{31},
  \qquad
  N_2=J_{02}+J_{32}
\]
generate \(G_q\). Their commutators are
\[
  [N_1,N_2]=0,
  \qquad
  [J_3,N_1]=\ii N_2,
  \qquad
  [J_3,N_2]=-\ii N_1.
\]
This is the Lie algebra of the Euclidean group \(ISO(2)\): \(J_3\) generates
rotations in a two-dimensional transverse plane, while \(N_1,N_2\) generate
commuting translations in that plane. The double cover is used on Hilbert
space.

\begin{center}
\begin{tikzpicture}[>=Stealth,scale=0.95]
  \draw[->,line width=0.7pt] (-3.0,0)--(3.2,0)
    node[right] {translation subgroup};
  \draw[->,line width=0.7pt] (0,-1.9)--(0,2.15)
    node[above] {rotation};
  \node[blue!70!black] at (0,1.35) {$J_3$};
  \draw[blue!70!black,very thick,->]
    (0.55,0.25) arc[start angle=25,end angle=325,radius=0.58];
  \node[green!50!black] at (-2.05,-0.35) {$N_1$};
  \node[green!50!black] at (2.05,-0.35) {$N_2$};
  \draw[green!50!black,very thick,->] (-1.55,-0.85)--(-0.45,-0.85);
  \draw[green!50!black,very thick,->] (0.45,-0.85)--(1.55,-0.85);
  \node[align=center] at (0,-1.45)
    {$ISO(2)=\mathbb R^2\rtimes SO(2)$};
\end{tikzpicture}
\end{center}

\section{Helicity Representations}

The helicity sectors considered in this volume are the massless
one-particle irreducible representations for which
\[
  N_1\ket{q,h}=0,
  \qquad
  N_2\ket{q,h}=0.
\]
The remaining little-group generator acts by
\[
  J_3\ket{q,h}=h\ket{q,h}.
\]
Equivalently, for
\[
  W(\alpha,\beta,\theta)
  =
  \exp(\ii\alpha N_1+\ii\beta N_2)\exp(\ii\theta J_3),
\]
the state transforms as
\[
  U(W(\alpha,\beta,\theta))\ket{q,h}
  =
  \ee^{\ii h\theta}\ket{q,h}.
\]
The double-cover rotation condition gives
\[
  h\in\frac12\mathbb Z.
\]

For general \(k\in\Sigma_0^+\), define generalized one-particle states by
\[
  \ket{k,h}=U(L(k))\ket{q,h},
\]
with covariant normalization
\[
  \langle k',h'|k,h\rangle
  =
  (2\pi)^3\,2k^0\,\delta^3(\vec k'-\vec k)\delta_{h'h}.
\]
For a Lorentz transformation \(\Lambda\), the Wigner little-group element is
\[
  W(\Lambda,k)=L(\Lambda k)^{-1}\Lambda L(k),
\]
and the helicity state transforms by
\[
  U(\Lambda)\ket{k,h}
  =
  \ee^{\ii h\theta(\Lambda,k)}\ket{\Lambda k,h},
\]
where \(\theta(\Lambda,k)\) is the \(SO(2)\) rotation angle of
\(W(\Lambda,k)\) in the helicity representation. Different choices of
\(L(k)\) change the phase convention for \(\ket{k,h}\) but leave invariant
matrix elements unchanged.

\begin{remark}
The massless little group also has irreducible unitary representations in
which the translation subgroup acts with nonzero orbit in the transverse
plane. These are continuous-spin representations. The helicity construction
above is the sector used for photons, gravitons, and the massless fields
developed in this volume.
\end{remark}

\section{Vector Polarization Representatives}

Consider helicity \(h=\pm1\). A vector polarization representative at the
reference momentum is
\[
  \epsilon_\mu^+(q)=\frac1{\sqrt2}(0,1,-\ii,0),
  \qquad
  \epsilon_\mu^-(q)=\frac1{\sqrt2}(0,1,\ii,0).
\]
It obeys
\[
  q^\mu\epsilon_\mu^\pm(q)=0,
  \qquad
  \epsilon_\mu^\pm(q)\epsilon^{\pm *\,\mu}(q)=1,
  \qquad
  \epsilon_\mu^\pm(q)\epsilon^{\mp *\,\mu}(q)=0.
\]
The rotation generated by \(J_3\) gives
\[
  W(0,0,\theta)_\mu{}^\nu\epsilon_\nu^\pm(q)
  =
  \ee^{\mp\ii\theta}\epsilon_\mu^\pm(q),
\]
which matches helicity \(\pm1\) after complex conjugation in the
vacuum-to-particle matrix element convention.

The null-translation part of the little group acts by a shift along \(q_\mu\).
For
\[
  W(\alpha,\beta,0)
  =
  \begin{pmatrix}
    1+\zeta&\alpha&\beta&-\zeta\\
    \alpha&1&0&-\alpha\\
    \beta&0&1&-\beta\\
    \zeta&\alpha&\beta&1-\zeta
  \end{pmatrix},
  \qquad
  \zeta=\frac{\alpha^2+\beta^2}{2},
\]
one finds
\[
  W(\alpha,\beta,0)_\mu{}^\nu\epsilon_\nu^\pm(q)
  =
  \epsilon_\mu^\pm(q)+c_\pm(\alpha,\beta)\,q_\mu,
\]
where \(c_\pm(\alpha,\beta)\) is a complex scalar depending on the
normalization of \(q\). Thus a covariant vector polarization for a helicity
\(\pm1\) massless particle is naturally a representative of the equivalence
class
\[
  \epsilon_\mu^\pm(k)\sim \epsilon_\mu^\pm(k)+c\,k_\mu .
\]
For general \(k\), choose
\[
  \epsilon_\mu^\pm(k)=L(k)_\mu{}^\nu\epsilon_\nu^\pm(q).
\]
Changing the standard transformation \(L(k)\) changes this representative by
the same equivalence relation and by the helicity phase.

\section{Field Strength Matrix Elements}

A vector potential \(A_\mu\) is a representative variable whose polarization
matrix element uses \(\epsilon_\mu^h(k)\). Its representative freedom is
\[
  A_\mu(x)\sim A_\mu(x)+\partial_\mu\xi(x),
\]
where \(\xi\) is a scalar gauge parameter. The gauge-invariant local tensor
constructed from \(A_\mu\) is
\[
  F_{\mu\nu}
  =
  \partial_\mu A_\nu-\partial_\nu A_\mu .
\]

For a one-photon state of helicity \(h=\pm1\), the field-strength matrix
element has the form
\[
  \langle k,h|\,\widehat F_{\mu\nu}(x)\,|\Omega\rangle
  =
  C_h\ee^{-\ii k\cdot x}
  \left(
    -\ii k_\mu\epsilon_\nu^{h*}(k)
    +\ii k_\nu\epsilon_\mu^{h*}(k)
  \right),
\]
where \(C_h\) is a normalization constant fixed by the field normalization.
The shift
\[
  \epsilon_\mu^h(k)\mapsto\epsilon_\mu^h(k)+c\,k_\mu
\]
leaves this expression unchanged because
\[
  k_\mu k_\nu-k_\nu k_\mu=0.
\]
The local photon observable at this stage is therefore the field strength,
while the potential is a representative used to write local equations,
actions, and perturbative Green functions.

\begin{center}
\begin{tikzpicture}[>=Stealth,scale=0.95]
  \draw[->,gray!70,line width=0.7pt] (-3.2,0)--(3.4,0)
    node[right] {representatives};
  \draw[->,gray!70,line width=0.7pt] (0,-1.6)--(0,1.9)
    node[above] {$k_\mu$ shift};
  \fill[blue!70!black] (-1.2,0) circle (2.4pt)
    node[below=0.15cm] {$\epsilon_\mu^h$};
  \fill[blue!70!black] (1.2,0.75) circle (2.4pt)
    node[above=0.12cm] {$\epsilon_\mu^h+c k_\mu$};
  \draw[blue!70!black,very thick,->] (-1.2,0)--(1.2,0.75);
  \node[align=center,fill=white,inner sep=1.5pt] at (-0.25,-1.08)
    {$k_\mu\epsilon_\nu^h-k_\nu\epsilon_\mu^h$\\is unchanged};
\end{tikzpicture}
\end{center}

\section{Polarization Sums}

For calculations with external helicity states it is useful to choose
representatives by imposing an auxiliary linear condition. Let \(n^\mu\) be a
fixed vector satisfying
\[
  n\cdot k\ne0 .
\]
Choose \(\epsilon_\mu^h(k;n)\) so that
\[
  k^\mu\epsilon_\mu^h(k;n)=0,
  \qquad
  n^\mu\epsilon_\mu^h(k;n)=0,
  \qquad
  \epsilon^h(k;n)\cdot\epsilon^{h'*}(k;n)=\delta_{hh'} .
\]
The sum over the two physical helicities is
\[
  \Pi_{\mu\nu}(k;n)
  =
  \sum_{h=\pm1}
  \epsilon_\mu^h(k;n)\epsilon_\nu^{h*}(k;n),
\]
with
\[
  \Pi_{\mu\nu}(k;n)
  =
  \eta_{\mu\nu}
  -
  \frac{k_\mu n_\nu+n_\mu k_\nu}{k\cdot n}
  +
  \frac{n^2 k_\mu k_\nu}{(k\cdot n)^2}.
\]
It satisfies
\[
  k^\mu\Pi_{\mu\nu}(k;n)=0,
  \qquad
  n^\mu\Pi_{\mu\nu}(k;n)=0.
\]
Changing \(n\) changes the representative polarization vectors by terms
proportional to \(k_\mu\). Therefore, if a quantity \(M^\mu(k)\) obeys
\[
  k_\mu M^\mu(k)=0,
\]
then
\[
  M^\mu(k)\Pi_{\mu\nu}(k;n)M^{\nu *}(k)
\]
is independent of the representative choice. This is the form in which
polarization sums will enter gauge-invariant scattering calculations.

The next chapter supplies the dynamics whose one-particle sector realizes
these helicity states: Maxwell theory, its constraints, and the controlled
introduction of gauge fixing.
